\addtocontents{toc}{\protect\Needspace{3\baselineskip}}
\originalchapter{均匀膨胀宇宙}\label{ch:cosmo}
\originalsection{尺度因子与红移}
在足够大的尺度上采用均匀、各向同性的近似. 共动坐标把随 Hubble 流运动的星系位置固定, 其实际距离随尺度因子 $a(t)$ 改变. 取今日 $a_0=1$, 则 $H=\dot a/a$. $v=HD$ 是同一宇宙时刻相邻共动观察者的距离变化, 大距离的退行率可以超过 $c$, 它不等于局部狭义相对论中的相对速度.

沿平直 FLRW 径向光线有 $\dd\chi=c\dd t/a(t)$. 相邻两个波峰连接同一对共动观察者, 积分端点微小变化给 $\delta t_o/a_o=\delta t_e/a_e$, 所以 $1+z=\lambda_o/\lambda_e=a_o/a_e$. 红移由整个传播期间的几何变化给出, 不需要把遥远星系与地球的比较强行视为单次 Lorentz boost.

\originalsection{Friedmann 动力学}
以下调用广义相对论卷已经建立的 FLRW 方程. $\rho$ 为包含全部组分的质量等效密度, $p$ 为压强:
\begin{align}
H^2&=\frac{8\pi G}3\rho-\frac{kc^2}{a^2},\label{eq:friedmann}\\
\frac{\ddot a}a&=-\frac{4\pi G}3\left(\rho+\frac{3p}{c^2}\right),\qquad
\dot\rho+3H\left(\rho+\frac p{c^2}\right)=0.\label{eq:conservation}
\end{align}
曲率项不表示一种真实物质. 非相互作用组分分别守恒, 常数状态参数 $w=p/(\rho c^2)$ 给 $\rho\propto a^{-3(1+w)}$. 物质 $w=0$, 辐射 $w=1/3$, 真空 $w=-1$. 后者虽然密度恒定, 共动体积内能量随体积增加, 由负压强做功关系补偿, 不应套用无压尘埃的能量直觉.

平直单组分宇宙若 $w>-1$, 由 $\dot a\propto a^{-(1+3w)/2}$ 积分得 $a\propto t^{2/[3(1+w)]}$. 物质期 $a\propto t^{2/3}$, 辐射期 $a\propto t^{1/2}$; 真空期须单独积分, 得 $a\propto\ee^{Ht}$. 总密度临界值为 $\rho_c=3H^2/(8\pi G)$, 定义 $\Omega_i=\rho_i/\rho_c$, 则 $\Omega_k=-kc^2/(a^2H^2)$ 且 $\sum\Omega_i+\Omega_k=1$.

\originalsection{宇宙时间与视界}
多组分模型写作
\begin{equation}\label{eq:Hz}
E(z)^2\equiv\frac{H(z)^2}{H_0^2}
=\Omega_{r0}(1+z)^4+\Omega_{m0}(1+z)^3
+\Omega_{k0}(1+z)^2+\Omega_{\Lambda0}.
\end{equation}
回望时间 $t_0-t(z)=H_0^{-1}\int_0^z\dd z'/[(1+z')E(z')]$. 年龄是上限改为无穷的积分, 须给出早期组分以保证模型在积分全域有效. 平直物质与真空模型有 $t_0=2\operatorname{arsinh}\sqrt{\Omega_\Lambda/\Omega_m}/(3H_0\sqrt{\Omega_\Lambda})$.

粒子视界 $D_p=a(t)\int_{t_{\rm start}}^t c\dd t'/a(t')$ 描述到今天能收到的最远因果传播; 事件视界 $D_e=a(t)\int_t^\infty c\dd t'/a(t')$ 描述现在发出的信号能到达的共动范围. 后者仅在积分收敛且模型持续有效时存在. Hubble 半径 $c/H$ 是局部膨胀尺度, 通常不等于前两者. 例如 $a\propto t^p$、$0<p<1$ 给 $D_p=ct/(1-p)$, 而持续指数膨胀给 $D_e=c/H$.
